\documentclass[a4paper,12pt]{article}
\usepackage[utf8]{inputenc}
\usepackage[T1]{fontenc}
\usepackage[french]{babel}
\usepackage{graphicx}
\usepackage{hyperref}
\usepackage{geometry}
\usepackage{titlesec}
\usepackage{enumitem}
\usepackage{xcolor}

\geometry{a4paper, margin=2.5cm}

\title{\textbf{Rapport de Travail - Projet Innovation Chess}}
\author{Développement Web et IA}
\date{\today}

\begin{document}

\maketitle

\tableofcontents
\newpage

\section{Introduction}

Ce rapport recense l'ensemble des fonctionnalités implémentées et des améliorations apportées au projet de site d'échecs avec entraînement par IA. Le projet combine une interface React moderne avec des fonctionnalités d'analyse de parties, d'entraînement tactique et positionnel, ainsi qu'une prédiction Elo utilisant l'apprentissage automatique.

\section{Fonctionnalités Implémentées}

\subsection{Game Rewinder (Rembobineur de Parties)}

\textbf{Description :}
Le Game Rewinder permet d'analyser une partie d'échecs en rembobinant les coups pour voir l'évolution de la position et la qualité des coups joués.

\textbf{Implémentations :}
\begin{itemize}
    \item Intégration dans un onglet "Analysis" séparé dans la page Game
    \item Calcul dynamique du FEN pour chaque position de rembobinage
    \item Synchronisation de l'échiquier avec l'index de rembobinage
    \item Mise à jour des pièces capturées en fonction de la position affichée
\end{itemize}

\textbf{Améliorations :}
\begin{itemize}
    \item Bouton de lecture automatique avec délai de 800ms entre les coups
    \item Indicateur visuel "Playing" avec animation pulse
    \item Raccourcis clavier : Flèches gauche/droite (navigation), Espace (avancer), Escape (quitter)
    \item Affichage de la catégorie du coup (best, good, inaccuracy, mistake, blunder)
    \item Évaluation numérique de chaque coup
\end{itemize}

% IMAGE 1: Screenshot du Game Rewinder avec les contrôles de rembobinage
\begin{figure}[h]
    \centering
    % TODO: Ajouter l'image ici
    % Screenshot du Game Rewinder montrant les contrôles de rembobinage, l'indicateur de coup actuel et l'analyse de la qualité du coup
    % Description: Interface du Game Rewinder avec le bouton de lecture automatique, les flèches de navigation et l'analyse du coup en cours
    \caption{Interface du Game Rewinder}
    \label{fig:rewinder}
\end{figure}

\subsection{Tactic Trainer (Entraînement Tactique)}

\textbf{Description :}
Le Tactic Trainer propose des puzzles tactiques interactifs pour améliorer la vision tactique du joueur.

\textbf{Implémentations :}
\begin{itemize}
    \item 6 puzzles niveau débutant avec positions FEN valides
    \item Thèmes tactiques variés : Fork, Pin, Discovered Attack, Skewer, Back Rank, Mate in 1
    \item Validation des solutions avec feedback immédiat
    \item Échiquier interactif pour résoudre les puzzles
\end{itemize}

\textbf{Améliorations :}
\begin{itemize}
    \item Système de progression : compteur de puzzles résolus
    \item Série de victoires (streak) avec indicateur visuel 🔥
    \item Effets sonores : playSuccess() pour les bonnes réponses, playError() pour les erreurs
    \item Statistiques en temps réel : puzzles résolus, précision, série actuelle
\end{itemize}

% IMAGE 2: Screenshot du Tactic Trainer avec un puzzle affiché
\begin{figure}[h]
    \centering
    % TODO: Ajouter l'image ici
    % Screenshot du Tactic Trainer montrant l'échiquier avec un puzzle tactique, le feedback et les statistiques
    % Description: Interface du Tactic Trainer avec un puzzle Fork, le feedback de l'IA Coach et les statistiques de progression
    \caption{Interface du Tactic Trainer}
    \label{fig:tactics}
\end{figure}

\subsection{Position Trainer (Entraînement Positionnel)}

\textbf{Description :}
Le Position Trainer permet de maîtriser des positions critiques d'ouvertures et de milieux de jeu.

\textbf{Implémentations :}
\begin{itemize}
    \item 6 positions niveau débutant avec FEN valides
    \item Thèmes positionnels : Italian Game, Pin, Center Control, Development, Piece Activity
    \item Identification du meilleur coup avec explications détaillées
    \item Feedback immédiat sur la qualité du coup joué
\end{itemize}

\textbf{Améliorations :}
\begin{itemize}
    \item Système de progression : positions résolues, précision
    \item Calcul de précision en temps réel
    \item Indicateur de thèmes maîtrisés
    \item Effets sonores adaptés aux réussites et erreurs
\end{itemize}

% IMAGE 3: Screenshot du Position Trainer avec une position affichée
\begin{figure}[h]
    \centering
    % TODO: Ajouter l'image ici
    % Screenshot du Position Trainer montrant l'échiquier avec une position positionnelle, le feedback et les statistiques
    % Description: Interface du Position Trainer avec une position Italian Game, l'explication du meilleur coup et les statistiques
    \caption{Interface du Position Trainer}
    \label{fig:position}
\end{figure}

\subsection{Elo Prediction avec ML}

\textbf{Description :}
Le système de prédiction Elo utilise l'apprentissage automatique pour estimer le niveau d'un joueur à partir de ses parties.

\textbf{Implémentations :}
\begin{itemize}
    \item Backend Python Flask avec API REST
    \item Intégration du modèle ML (modele\_elo\_optimise\_final.pkl)
    \item Extraction de caractéristiques des parties PGN
    \item Communication avec le frontend React via API
\end{itemize}

\textbf{Améliorations :}
\begin{itemize}
    \item Méthode de fallback si le modèle ML ne peut pas être chargé
    \item Prédiction basée sur des heuristiques (nombre de coups, captures, échecs)
    \item Gestion robuste des erreurs
    \item Indicateur de la méthode utilisée (ML ou fallback)
\end{itemize}

% IMAGE 4: Screenshot de la page Elo Prediction avec les résultats
\begin{figure}[h]
    \centering
    % TODO: Ajouter l'image ici
    % Screenshot de la page Elo Prediction montrant l'upload de PGN, les résultats de prédiction et les métriques
    % Description: Interface de prédiction Elo avec les résultats d'analyse, l'Elo prédit et les métriques détaillées
    \caption{Interface de Prédiction Elo}
    \label{fig:elo}
\end{figure}

\subsection{Restructuration de l'AI Coach}

\textbf{Description :}
L'AI Coach a été restructuré en pages séparées pour améliorer l'organisation et l'utilisation de l'espace.

\textbf{Implémentations :}
\begin{itemize}
    \item Page d'accueil Coach avec liens vers les sections
    \item Pages séparées : CoachElo, CoachTactics, CoachPosition, CoachThemes
    \item Menu déroulant dans la navigation pour accès rapide
    \item Routage mis à jour dans App.tsx
\end{itemize}

\textbf{Avantages :}
\begin{itemize}
    \item Meilleure organisation du code
    \item Réduction du scrolling nécessaire
    \item Navigation plus intuitive
    \item Chargement plus rapide des pages individuelles
\end{itemize}

% IMAGE 5: Screenshot de la page d'accueil Coach avec les cartes de navigation
\begin{figure}[h]
    \centering
    % TODO: Ajouter l'image ici
    % Screenshot de la page Coach montrant les 4 cartes de navigation vers Elo Prediction, Tactic Trainer, Position Trainer et Theme Analyzer
    % Description: Page d'accueil de l'AI Coach avec les cartes de navigation colorées et les descriptions
    \caption{Page d'accueil de l'AI Coach}
    \label{fig:coach}
\end{figure}

\subsection{Optimisation de l'Espace et de l'UI}

\textbf{Description :}
L'interface utilisateur a été optimisée pour réduire le scrolling et améliorer l'expérience utilisateur.

\textbf{Implémentations :}
\begin{itemize}
    \item Game Rewinder déplacé dans un onglet Analysis séparé
    \item Sidebar optimisée pour afficher les informations essentielles
    \item Utilisation efficace de l'espace horizontal
    \item Design cohérent avec Tailwind CSS
\end{itemize}

\textbf{Améliorations :}
\begin{itemize}
    \item Feedback visuel pendant le rembobinage
    \item Animations fluides pour les transitions
    \item Indicateurs de progression clairs
    \item Interface responsive et moderne
\end{itemize}

\section{Architecture Technique}

\subsection{Frontend (React + TypeScript)}
\begin{itemize}
    \item React 18 avec TypeScript
    \item Vite comme outil de build
    \item Chess.js pour la logique d'échecs
    \item Tailwind CSS pour le styling
    \item React Router pour la navigation
\end{itemize}

\subsection{Backend ML (Python Flask)}
\begin{itemize}
    \item Flask pour l'API REST
    \item Flask-CORS pour la communication cross-origin
    \item NumPy pour le traitement des données
    \item Scikit-learn pour le modèle ML
    \item Pickle pour charger le modèle entraîné
\end{itemize}

\subsection{Base de Données}
\begin{itemize}
    \item Supabase pour l'authentification
    \item Supabase pour le stockage des données
    \item Configuration via variables d'environnement
\end{itemize}

\section{Améliorations Possibles}

\subsection{Fonctionnalités Additionnelles}

\subsubsection{Système de Progression Avancé}
\begin{itemize}
    \item Système XP et niveaux pour les joueurs
    \item Classement global avec leaderboard
    \item Badges et achievements
    \item Défis quotidiens et hebdomadaires
\end{itemize}

\subsubsection{Contenu Éducatif Étendu}
\begin{itemize}
    \item Puzzles intermédiaires et avancés
    \item Leçons théoriques interactives
    \item Études de finales
    \item Analyseur d'ouvertures
    \item Exercices de calcul
\end{itemize}

\subsubsection{Mode de Jeu Amélioré}
\begin{itemize}
    \item Mode Blitz avec timer pour les puzzles
    \item Parties personnalisées
    \item Analyse en temps réel pendant les parties
    \item Suggestions de coups avec Stockfish
\end{itemize}

\subsection{Améliorations UX/UI}

\subsubsection{Interface Utilisateur}
\begin{itemize}
    \item Mode sombre/clair
    \item Sons personnalisables
    \item Animations plus fluides entre les pages
    \item Notifications de progression
    \item Interface mobile responsive
\end{itemize}

\subsubsection{Accessibilité}
\begin{itemize}
    \item Support des lecteurs d'écran
    \item Navigation au clavier complète
    \item Contrastes améliorés
    \item Police ajustable
\end{itemize}

\subsection{Améliorations Techniques}

\subsubsection{Performance}
\begin{itemize}
    \item Code-splitting avec dynamic import()
    \item Lazy loading des composants
    \item Optimisation des bundles
    \item Mise en cache des puzzles
    \item Service Worker pour offline (PWA)
\end{itemize}

\subsubsection{Sécurité}
\begin{itemize}
    \item Validation des entrées côté serveur
    \item Rate limiting sur l'API
    \item Protection CSRF
    \item Chiffrement des données sensibles
\end{itemize}

\subsubsection{Tests}
\begin{itemize}
    \item Tests unitaires avec Jest
    \item Tests E2E avec Playwright
    \item Tests d'intégration
    \item Couverture de code
\end{itemize}

\subsection{Améliorations du Modèle ML}

\subsubsection{Entraînement du Modèle}
\begin{itemize}
    \item Réentraînement avec plus de données
    \item Feature engineering avancé
    \item Hyperparameter tuning
    \item Validation croisée
    \item A/B testing des prédictions
\end{itemize}

\subsubsection{Fonctionnalités ML}
\begin{itemize}
    \item Prédiction du style de jeu
    \item Analyse des faiblesses
    \item Recommandations d'entraînement personnalisées
    \item Comparaison avec des joueurs de même niveau
    \item Analyse temporelle de la progression
\end{itemize}

\subsection{Intégrations Externes}

\subsubsection{Plateformes d'Échecs}
\begin{itemize}
    \item Import automatique depuis Chess.com
    \item Import automatique depuis Lichess
    \item Synchronisation des parties
    \item Analyse des parties en ligne
\end{itemize}

\subsubsection{Outils d'Analyse}
\begin{itemize}
    \item Intégration de Stockfish pour analyse approfondie
    \item Tableaux d'ouvertures
    \item Base de données de finales
    \item Annotation automatique des parties
\end{itemize}

\section{Conclusion}

Ce projet a permis de développer une plateforme complète d'entraînement aux échecs combinant interface moderne, fonctionnalités interactives et intelligence artificielle. Les implémentations réalisées offrent une base solide pour un outil d'apprentissage efficace et engageant.

Les améliorations possibles identifiées permettent d'envisager une évolution continue de la plateforme vers un produit plus complet et plus performant, capable de rivaliser avec les meilleures plateformes d'entraînement aux échecs existantes.

\end{document}
